Every value of a Python program is an object:
an integer, a dictionary, a function, a book.

\begin{defi}\label{def.object}
 An \emph{object} has an \emph{identity}, a \emph{type} and a \emph{value}
 \citep[Section 3.1]{PSF26lan}.
 The identity and the type are fixed when the object is created.
 The object is \emph{mutable} if its value can change afterwards,
 and \emph{immutable} otherwise.
 The expression \codehl{a is b} compares identities,
 and \codehl{a == b} compares values.
\end{defi}

Integers, floats, strings and tuples are immutable;
lists and dictionaries are mutable.
A \emph{name} refers to an object.
An assignment \emph{binds} a name to an object,
as the statements \codehl{def} and \codehl{class} do,
and it copies nothing:
after \codehl{b = a}, the names \codehl{a} and \codehl{b} refer to one object.
A change made to a mutable object through one name is read through every other,
and the parameters of a function are names of the objects that the caller passed.
In Listing \ref{listing.namesAndObjects}
a side of the book of Section \ref{sec.foldImplementation},
which is a dictionary from price to size,
is written through one name and read through another;
a second dictionary with the same entries is equal to it, and at no time identical.

\begin{lstlisting}[caption={Two names of one dictionary, and two equal dictionaries.}, label=listing.namesAndObjects]
bids = {99: 100, 98: 60}
same = bids
bids[99] = 40
assert same is bids and same[99] == 40

copy = dict(bids)
assert copy == bids and copy is not bids
bids[99] = 10
assert copy[99] == 40 and copy != bids
\end{lstlisting}

An immutable object needs no such care, since no name can change it;
Section \ref{sec.dataClasses} makes the messages of the stream immutable for that reason.
